
The gradient flow in field theory has proven to be a very useful
concept.  Originally introduced as a method to regularize divergences
and ultraviolet fluctuations in lattice
calculations\,\cite{Narayanan:2006rf,Luscher:2009eq}, it has mainly
gained traction through its use for scale-setting in lattice
\qcd{}\,\cite{Luscher:2010iy,Luscher:2013vga,Borsanyi:2012zs} and has
meanwhile become an indispensable tool in many practical calculations in
this field (see,
e.g.\ \citeres{Luscher:2013cpa,Sommer:2014mea,Ramos:2015dla,Aoki:2019cca}).
While first focused on \qcd{}, generalizations of the gradient-flow
formalism (\gff) yield a much wider range of applications, for example
the study of dualities in field
theory\,\cite{Aoki:2015dla,Aoki:2016ohw,Aoki:2017uce,Aoki:2017bru,
  Aoki:2018dmc,Aoki:2019bfb}.

In \qcd{}, the formulation as a five-dimensional field theory has
been presented in \citere{Luscher:2011bx}, where the additional
dimension is associated with the so-called flow time. The form of the
fundamental theory (i.e., actual \qcd{}) serves as the boundary
condition at vanishing flow time. It has been proven that, when
expressed in terms of renormalized parameters and fields, composite
operators at positive flow time are finite\,\cite{Luscher:2011bx}. This
property and the realization that the gradient flow can be used as a
matching scheme between lattice and perturbative calculations shed light
on a whole new set of applications that are currently being explored.

One of the first possible cross-fertilizations among perturbative and
lattice \qcd{} is given by the definition of a new scheme for the strong
coupling, defined by the gluon condensate (\qcd{} action density) at
finite flow time\,\cite{Luscher:2010iy}.  Its perturbative relation to
the \msbar{} coupling is known through three
loops\,\cite{Harlander:2016vzb}.\footnote{Unless stated otherwise, we
  refer to perturbative calculations as performed at infinite volume in
  this paper. The inclusion of finite-volume effects requires different
  techniques, such as Numerical Stochastic Perturbation Theory, see
  \citere{DallaBrida:2017tru}.}  Another cross-boundary application is
the gradient-flow definition of the energy-momentum tensor, which uses
the small-flow-time expansion in order to express it in terms of
well-defined composite operators at positive flow time and
perturbatively accessible coefficient
functions\,\cite{Suzuki:2013gza,Makino:2014taa,
  Harlander:2018zpi,Iritani:2018idk}. Yet another example is the
proposal to relate Euclidean quasi \pdf{}s\footnote{\pdf=parton
  distribution function.} on the lattice to perturbative light-front
\pdf{}s by using the gradient
flow\,\cite{Monahan:2016bvm,Monahan:2017hpu}.

All these applications rely on input from the lattice as well as from
perturbative calculations, both at finite flow time. It was found that
higher order corrections in perturbation theory are crucial for reducing
the perturbative truncation error as estimated by the variation of the
renormalization scale\,\cite{Harlander:2016vzb,Harlander:2018zpi}. Such
corrections were also found to significantly stabilize the required
extrapolations of the corresponding lattice
results\,\cite{Iritani:2018idk}.

Despite the fact that loop integrals in the \gff\ involve additional
exponential factors as well as integrations over flow-time variables,
many important techniques for regular perturbative calculations retain
their usefulness, albeit in slightly generalized form. It is the goal of
this paper to provide the basis for further perturbative calculations
within the \gff, and thus to contribute to the further exploration of
the capabilities of this approach. \sct{sec:pert} reviews the \qcd{}
gradient flow in perturbation theory, recapitulating results of
\citeres{Luscher:2010iy,Luscher:2011bx,Luscher:2013cpa}. The various
stages of automation in a perturbative multi-loop calculation of
correlation functions are described in \sct{sec:implementation}: the
generation of Feynman diagrams and insertion and evaluation of Feynman
rules, followed by a reduction of the loop integrals to a set of master
integrals, and finally the numerical evaluation of the latter. As an
application, \sct{sec:observables} presents the three-loop results for
the gluon condensate, obtained before in \citere{Harlander:2016vzb}, the
quark condensate, as well as the ``ringed'' quark field renormalization
constant, which are new results.  These quantities allow one to define a
precision gradient-flow coupling and gradient-flow mass scheme and their
matching to \msbar{} as shown in \sct{sec:couplingandmass}. Our
conclusions are presented in \sct{sec:conclusions}.
